\documentclass[10pt,a4paper]{amsart}
\usepackage[margin=19mm]{geometry}
\usepackage{amsmath,amssymb,mathtools}
\usepackage[scr=rsfs,cal=euler]{mathalfa}
\usepackage{xfrac}
\usepackage[unicode,hidelinks,bookmarks=false]{hyperref}
\newcommand{\ie}{i.e.\ }
\newcommand{\cf}{cf.\ }
\newcommand{\resp}{respectively\ }
\newcommand{\Subsection}{\subsection}
\providecommand{\nobreakdash}{\nobreak\hbox{-}}
\setlength{\emergencystretch}{2em}
\multlinegap0pt
\usepackage{amssymb}
\usepackage{mathtools}
\newtheorem{proposition}{Proposition}
\newcommand{\W}{\mathcal{W}}
\title{A note on a local entropy condition in the Wasserstein space}
\author{Chang Jun Im, Jeong Min Jeon, Byeong U. Park}
\date{Source: arXiv:2609.08403v1 (CC BY 4.0)}
\begin{document}
\maketitle

\noindent\emph{Source and licence.} A note on a local entropy condition in the Wasserstein space. Version arXiv:2609.08403v1 (CC BY 4.0), distributed by arXiv under the Creative Commons Attribution 4.0 International licence, http://creativecommons.org/licenses/by/4.0/, as recorded in the arXiv licence metadata for this exact version and shown on the abstract page https://arxiv.org/abs/2609.08403v1. Full licence notice: \texttt{licenses/arxiv-2609.08403-CC-BY-4.0.txt}.\par\smallskip

\noindent\textit{Editorial scope.} Counted unit: the article's proposition thm:second (source0.tex lines 229--232). The article's complete original proof of it is delivered with it (source0.tex lines 234--330, one \texttt{proof} environment of 97 lines). No mathematics is written, repaired or re-derived: the statement and the proof are the article's own lines, in the article's own notation, and no retained line is substituted or re-flowed.  Retained uncounted as the source-local context the proof needs: lines 133--135.  Nothing was omitted from any retained span, so the package carries no omission record.  No retained line contains a citation, so the wrapper carries no bibliography and no bibliography entry.  No retained line contains a figure environment, an image-inclusion command or drawing code, so no image is delivered and the package declares no asset dependency.  Editorial formatting only, in addition: an AMS \texttt{amsart} preamble that restates the article's own declarations and macros used by the retained lines, namely the environments \texttt{proposition} on separate counters, together with the standard packages those lines need.  The retained environments print in this document as Proposition 1 in order of appearance, whereas the article prints the same items with the numbering of its own full document.  The complete native source of the article as distributed in the arXiv e-print for this exact version is bundled unchanged as \texttt{sources/209712/source0.tex}.
\begin{proposition} \label{thm:first}
    Let $F \in \mathcal{W}$ be a distribution function that is strictly increasing and absolutely continuous on $[0,1]$. Suppose that there exist constants \(0<c_0\le c_1<\infty\) such that $c_0\le F'(u)\le c_1$ for almost everywhere $u\in(0,1)$. Then, $H(F,\delta) \to \infty$ as $\delta \to 0$.
\end{proposition}

\begin{proposition} \label{thm:second}
    Let \(c_0\) and \(c_1\) satisfy \(L<c_0<c_1<U\), and let
\(F\in\W_{L,U}\) be a distribution function such that $c_0|u-v|\le |F(u)-F(v)|\le c_1|u-v|$ for all $u,v\in[0,1]$. Then, $H_{L, U}(F,\delta) \to \infty$ as $\delta \to 0$.
\end{proposition}

\begin{proof}
    Fix a small $\delta > 0$. Consider sequences $\{a_l: l \ge 0\}$ and $\{b_l: l \ge 1\}$ defined as
    \begin{align*}
        a_l &\equiv a_l(\delta):= \left( \frac{27\delta^2}{16U} \right)^{1/3} \times l, \\
        b_l &\equiv b_l(\delta):= \frac{U a_{l} - L a_{l-1} + F(a_{l-1}) - F(a_{l})}{U-L}.
    \end{align*}
    Define $K \equiv K(\delta):= \max \{l \in \mathbb{N} : a_l \leq 1 \}$. For $l \in \{1, \ldots, K\}$, it follows that
    \begin{align*}
        L (a_{l} - a_{l-1}) < c_0 (a_{l} - a_{l-1}) \le F(a_{l}) - F(a_{l-1}) \le c_1 (a_{l} - a_{l-1}) < U (a_{l} - a_{l-1}),
    \end{align*}
    which implies that $a_{l-1} < b_l < a_l$. For $l \in \{1, \ldots, K\}$, define $F_l \equiv F_l(\cdot,\delta)\in\mathcal{W}_{L,U}$ as
    \begin{align*}
        F_l (u) :=
        \begin{cases}
            0, & \text{if}~u < 0, \\
            F(u), & \text{if}~u\in[0,a_{l-1})\cup[a_l, \infty),\\
            L(u-a_{l-1}) + F(a_{l-1}), & \text{if}~u \in [a_{l-1}, b_l), \\
            U(u-a_{l}) + F(a_{l}), & \text{if}~u \in [b_l, a_l). %, \\
        \end{cases}
    \end{align*}
    The quantile function of $F_l$ is given by
    \begin{align*}
        F_l^{-1} (t) =
        \begin{cases}
            \frac{t-F(a_{l-1})}{L} + a_{l-1}, & \text{if}~t \in (F(a_{l-1}), F_l(b_l)], \\
            \frac{t-F(a_{l})}{U} + a_{l}, & \text{if}~t \in (F_l(b_{l}), F(a_l)], \\
            F^{-1}(t), & \text{if}~t \notin (F(a_{l-1}), F(a_l)].
        \end{cases}
    \end{align*}
    A direct calculation shows that
    \begin{align*}
        \frac{t-F(a_{l-1})}{U} + a_{l-1} < \frac{t-F(a_{l-1})}{c_1}+ a_{l-1} &\le F^{-1} (t) \le \frac{t-F(a_{l})}{c_1} + a_{l} \le \frac{t-F(a_{l})}{U} + a_{l}, \\
        \frac{t-F(a_{l-1})}{U} + a_{l-1} &\le F_l^{-1} (t) \le \frac{t-F(a_{l})}{U} + a_{l}
    \end{align*}
    for all $t\in(F(a_{l-1}),F(a_l)]$.
    This implies that
    \begin{align}\label{diff bound}
    \begin{split}
    |F^{-1}(t)-F^{-1}_l(t)| & < \left|\frac{t-F(a_{l})}{U} + a_{l}-\left(\frac{t-F(a_{l-1})}{U} + a_{l-1}\right) \right| \\
    & = \left|a_{l} - a_{l-1} - \frac{F(a_l) -F(a_{l-1})}{U}\right|.
    \end{split}
    \end{align}
    It follows that
    \begin{align*}
        d_{\mathcal{W}} \left(F, F_l \right)^2 & = \int_0^1 (F^{-1}(t)- F_l^{-1}(t))^2 \mathrm{d}t
        = \int_{F(a_{l-1})}^{F(a_l)} \left(F^{-1}(t) - F_l^{-1}(t) \right)^2 \mathrm{d}t \\
        &< \int_{F(a_{l-1})}^{F(a_l)} \left(a_{l} - a_{l-1} - \frac{F(a_l) -F(a_{l-1})}{U}\right)^2 \mathrm{d}t \\
        & = \left(a_{l} - a_{l-1} - \frac{F(a_l) -F(a_{l-1})}{U}\right)^2 \left( F(a_l) -  F(a_{l-1}) \right) \\
        & = \frac{1}{2U^2} \left(U (a_{l} - a_{l-1}) - (F(a_l) -F(a_{l-1})) \right)^2 \left( 2( F(a_l) -  F(a_{l-1}) ) \right) \\
        & \le \frac{1}{2U^2} \left( \frac{2U(a_{l} - a_{l-1})}{3} \right)^3 = \frac{4U}{27} (a_{l} - a_{l-1})^3 = \frac{\delta^2}{4},
    \end{align*}
    where the first inequality follows from \eqref{diff bound}, and the last follows from the arithmetic--geometric mean inequality applied to three nonnegative numbers. Now, define
    \begin{align*}
        \epsilon_0 := \min \left\{ \left(\frac{9L^3}{128U}\right)^{1/2} \cdot \min \left\{ \frac{1}{L} - \frac{1}{c_0} , \frac{1}{c_1}- \frac{1}{U} \right\}, \frac{1}{2} \right\}
    \end{align*}
    and let $\epsilon \in \left(0, \epsilon_0 \right)$. Since
    \begin{align*}
        F^{-1} (t) \le \frac{t-F(a_{l-1})}{c_0} + a_{l-1} < F_l^{-1} (t), \quad &\text{if}~t \in (F(a_{l-1}), F_l(b_l)], \\
        F^{-1} (t) \le \frac{t-F(a_{l})}{c_1} + a_{l} \le F_l^{-1} (t), \quad &\text{if}~t \in (F_l(b_l), F_l(a_l)],
    \end{align*}
    we have
    \begin{align}\label{some bounds}
    \begin{split}
    F_l^{-1}(t)-F^{-1}(t)\ge F_l^{-1}(t)-\left(\frac{t-F(a_{l-1})}{c_0} + a_{l-1}\right), \quad &\text{if}~t \in (F(a_{l-1}), F_l(b_l)], \\
    F_l^{-1}(t)-F^{-1}(t)\ge F_l^{-1}(t)-\left(\frac{t-F(a_{l})}{c_1} + a_{l}\right), \quad &\text{if}~t \in (F_l(b_l), F_l(a_l)].
    \end{split}
    \end{align}
    It follows that
    \begin{align}\begin{split} \label{eq:7.33_2}
        & \int_{F(a_{l-1})}^{F(a_l)} \left(F^{-1}(t) - F_l^{-1}(t) \right)^2 \mathrm{d}t \\
        & = \int_{F(a_{l-1})}^{F_l(b_l)} \left(F^{-1}(t) - F_l^{-1}(t) \right)^2 \mathrm{d}t + \int_{F_l(b_{l})}^{F(a_l)} \left(F^{-1}(t) - F_l^{-1}(t) \right)^2 \mathrm{d}t \\
        & \ge \int_{F(a_{l-1})}^{F_l(b_l)} \left( \frac{1}{L} - \frac{1}{c_0} \right)^2 \left( t-F(a_{l-1}) \right)^2 \mathrm{d}t + \int_{F_l(b_{l})}^{F(a_l)} \left( \frac{1}{c_1} - \frac{1}{U} \right)^2 \left( t-F(a_{l}) \right)^2 \mathrm{d}t \\
        & = \frac{1}{3} \left( \frac{1}{L} - \frac{1}{c_0} \right)^2 \left( F_l(b_l)-F(a_{l-1}) \right)^3 + \frac{1}{3} \left( \frac{1}{c_1} - \frac{1}{U} \right)^2 \left( F(a_l)-F_l(b_l) \right)^3 \\
        & \ge \frac{1}{3} \min \left\{ \left( \frac{1}{L} - \frac{1}{c_0} \right)^2, \left( \frac{1}{c_1} - \frac{1}{U} \right)^2 \right\} \left( \left( F_l(b_l)-F(a_{l-1}) \right)^3 + \left( F(a_l)-F_l(b_l) \right)^3 \right) \\
        & \ge \frac{1}{12} \min \left\{ \left( \frac{1}{L} - \frac{1}{c_0} \right)^2, \left( \frac{1}{c_1} - \frac{1}{U} \right)^2 \right\} \left( F(a_l)-F(a_{l-1}) \right)^3 \\
        & > \frac{L^3}{12} \min \left\{ \left( \frac{1}{L} - \frac{1}{c_0} \right)^2, \left( \frac{1}{c_1} - \frac{1}{U} \right)^2 \right\} \left( a_l-a_{l-1} \right)^3 \\
        & = \frac{9L^3}{64U} \min \left\{ \left( \frac{1}{L} - \frac{1}{c_0} \right)^2, \left( \frac{1}{c_1} - \frac{1}{U} \right)^2 \right\} \delta^2 \ge 2 \delta^2 \epsilon_0^2 > 2 \delta^2 \epsilon^2,
    \end{split} \end{align}
    where the first inequality follows from \eqref{some bounds} and the third inequality follows from the fact that $4(a^3+b^3)\ge(a+b)^3$ for all $a,b>0$.
    By \eqref{eq:7.33_2}, for all $1\leq l<m\leq K$, we have
    \begin{align}\begin{split} \label{eq:7.33_3}
        d_{\mathcal{W}} \left( F_m,  F_l \right)^2
        & = \int_{F(a_{l-1})}^{F(a_l)} \left( F^{-1}_m(t) - F_l^{-1}(t) \right)^2 \mathrm{d}t + \int_{F(a_{m-1})}^{F(a_m)} \left( F^{-1}_m(t) - F_l^{-1}(t) \right)^2 \mathrm{d}t \\
        & = \int_{F(a_{l-1})}^{F(a_l)} \left( F^{-1}(t) - F_l^{-1}(t) \right)^2 \mathrm{d}t + \int_{F(a_{m-1})}^{F(a_m)} \left( F^{-1}(t) - F_m^{-1}(t) \right)^2 \mathrm{d}t > 4 \delta^2 \epsilon^2.
    \end{split} \end{align}
    Arguing as in the proof of Proposition \ref{thm:first}, we get
    \begin{align*}
        K=N\Big(\delta\epsilon,\,\bigcup_{l=1}^K B_{\mathcal{W}_{L, U}}(F_l, \delta \epsilon),\,d_{\mathcal{W}}\Big) \le N(\delta\epsilon/2,B_{\mathcal{W}_{L, U}}(F, \delta),d_{\mathcal{W}}).
    \end{align*}
    This implies that
    \begin{align*}
        \int_{0}^{1} \sqrt{1+\log N(\delta \epsilon, B_{\mathcal{W}_{L, U}}(F, \delta), d_{\mathcal{W}})} \, \mathrm{d} \epsilon
        &= \frac{1}{2} \int_{0}^2 \sqrt{1+\log N(\delta \epsilon/2, B_{\mathcal{W}_{L, U}}(F, \delta), d_{\mathcal{W}})} \, \mathrm{d} \epsilon \\
        &\ge \frac{\epsilon_0}{2}\sqrt{1+\log\Big\lfloor\Big(\frac{16U}{27\delta^2}\Big)^{1/3}\Big\rfloor},
    \end{align*}
    where the inequality follows from the definition of $K$. This completes the proof of the proposition.
\end{proof}

\end{document}
